% Compile with XeLaTeX twice. Revised research manuscript, 2026-10-13.
\documentclass[a4paper,11pt]{article}
\usepackage{fontspec,xeCJK}
\setmainfont{Times New Roman}
\setCJKmainfont[BoldFont=SimHei,ItalicFont=KaiTi]{SimSun}
\setCJKsansfont{SimHei}
\setCJKmonofont{FangSong}
\usepackage[margin=24mm,headheight=15pt]{geometry}
\usepackage{amsmath,amssymb,amsthm,mathtools,booktabs,array}
\usepackage{hyperref,fancyhdr}
\hypersetup{hidelinks,pdftitle={孙智伟一个加权二项式和整除猜想的证明},pdfauthor={}}
\pagestyle{fancy}
\fancyhf{}
\fancyhead[L]{\small 加权二项式和的整除性}
\fancyhead[R]{}
\fancyfoot[C]{\thepage}
\renewcommand{\headrulewidth}{0.3pt}
\setlength{\parindent}{2em}
\setlength{\parskip}{0.15em}
\setlength{\emergencystretch}{2em}
\linespread{1.10}
\allowdisplaybreaks[2]
\numberwithin{equation}{section}
\renewcommand{\abstractname}{摘要}
\renewcommand{\refname}{参考文献}
\renewcommand{\proofname}{证明}
\newtheoremstyle{chinese}{6pt}{6pt}{\normalfont}{} {\bfseries}{}{0.5em}{}
\theoremstyle{chinese}
\newtheorem{theorem}{定理}
\newtheorem{lemma}{引理}[section]
\newtheorem{proposition}[lemma]{命题}
\newcommand{\Z}{\mathbb Z}
\newcommand{\Q}{\mathbb Q}
\newcommand{\diag}{\Delta}
\newcommand{\J}{\mathcal J}
\title{\textbf{孙智伟一个加权二项式和整除猜想的证明}}
\author{}
\date{}
\begin{document}
\maketitle
\begin{abstract}
设 $a_n=\sum_{i,j=0}^{n}\binom ni^2\binom nj^2\binom{i+j}{i}2^{i+j}3^{2n-i-j}$。
证明对每个正整数 $n$，$n$ 整除 $2\sum_{k=0}^{n-1}(3k+2)a_k$，
从而证明孙智伟猜想 4.3(ii) 中的整性断言。
利用对角系数表示和多项式微分恒等式建立统一分母界，再证明二进与三进整性。
三进整性的证明用到奇次幂展开和有限原函数的余项估计。
\end{abstract}
\noindent\textbf{关键词：}二项式和；整除性；对角系数；多项式证书；局部整性

\noindent\textbf{MR(2020)分类号：}11B65，11A07，05A19\qquad
\textbf{中图分类号：}O156.1

\bigskip
\begin{center}
\textbf{Proof of a divisibility conjecture of Zhi-Wei Sun\\
on a weighted binomial sum}\\[6pt]

\end{center}
\noindent\textbf{Abstract.}
Let $a_n=\sum_{i,j=0}^{n}\binom ni^2\binom nj^2\binom{i+j}{i}2^{i+j}3^{2n-i-j}$.
We prove that $n$ divides $2\sum_{k=0}^{n-1}(3k+2)a_k$ for every positive integer $n$.
This proves the integrality assertion in Sun's Conjecture 4.3(ii).
A diagonal coefficient representation and an explicit differential certificate
give a uniform denominator bound. Integrality at the primes 2 and 3 then removes
the remaining possible denominator factors. For the prime 3, we express the generating functions as finite sums of odd powers
and estimate the remainders of truncated primitives.

\medskip
\noindent\textbf{Keywords:} binomial sum; divisibility; diagonal coefficient;
polynomial certificate; local integrality.

\newpage
\section{引言}
孙智伟\cite{Sun2012}研究了 Apéry 多项式
\[
\mathcal A_n(x)=\sum_{k=0}^n\binom nk^2\binom{n+k}{k}^2x^k
\]
及其加权和，证明了
$n^{-1}\sum_{k=0}^{n-1}(2k+1)\mathcal A_k(x)\in\Z[x]$。
Guo\cite{Guo2016}通过 Delannoy 多项式平方的求和公式，
证明了若干涉及偶次幂的有理和在整数点取整数值。

对非负整数 $d,n$，孙智伟\cite{Sun2026}引入了多项式族
\begin{equation}\label{eq:Sun-polynomial}
S_n^{(d)}(x)=\sum_{i,j=0}^n\binom ni^d\binom nj^d\binom{i+j}{i}x^{i+j}
\qquad(d,n\ge0),
\end{equation}
并提出了相关的整性与素数幂同余猜想。
本文研究 $d=2$、$x=2/3$ 时的整数序列
$a_n=9^nS_n^{(2)}(2/3)$，证明如下结果。
\begin{theorem}\label{thm:main}
对每个非负整数 $n$，定义
\begin{equation}\label{eq:sequence}
a_n=\sum_{i=0}^{n}\sum_{j=0}^{n}
\binom ni^2\binom nj^2\binom{i+j}{i}
2^{i+j}3^{2n-i-j}.
\end{equation}
则对每个正整数 $n$，都有
\begin{equation}\label{eq:goal}
n\ \bigg|\ 2\sum_{k=0}^{n-1}(3k+2)a_k.
\end{equation}
\end{theorem}

定理~\ref{thm:main}~证明了文献\cite[Conjecture 4.3(ii)]{Sun2026}中的第一条整性断言。

证明利用二项式变换和多项式微分恒等式建立统一分母界，再证明二进与三进整性。
二进整性由模 $2^s$ 的周期性得到；三进部分将相应的生成函数写成
$(1-16x)^{\alpha/2}$ 的有限和，其中 $\alpha$ 为奇数，继而估计有限原函数的余项。
求和问题中的证书方法可参见 Zeilberger\cite{Zeilberger1991}。

\subsection*{记号与形式系数约定}
对每个正整数 $n$，记
\begin{equation}\label{eq:T}
T_n=\frac2n\sum_{k=0}^{n-1}(3k+2)a_k\in\Q.
\end{equation}
对素数 $p$，$v_p(q)$ 表示非零有理数 $q$ 中 $p$ 的指数，并约定 $v_p(0)=+\infty$。
记 $\Z_{(p)}=\{q\in\Q:v_p(q)\ge0\}$。
我们依次证明
\begin{equation}\label{eq:threeclaims}
24T_n\in\Z,\qquad v_2(T_n)\ge0,\qquad v_3(T_n)\ge0.
\end{equation}
由这三式可得 $T_n\in\Z$。

以下生成函数均视为形式幂级数，$[x^j]F$ 表示 $F$ 的第 $j$ 个系数。
对非负整数 $k$，约定 $j>k$ 时 $\binom{k}{j}=0$。

\section{统一分母界}
\subsection{对角系数与二项式变换}
令 $\diag_r F=[u^rv^rz^r]F$，并定义
\begin{align}
\Psi&=(1+2u)(1+2v)(1+3u+z)(1+z+3vz),\\
\Phi&=\Psi-uvz.\label{eq:Phi}
\end{align}
先取 $u,v$ 的系数，有
\begin{equation}
\begin{split}
[u^kv^k]\Psi^k
=\sum_{i,j=0}^k\binom ki^2\binom kj^2
2^{i+j}3^{2k-i-j}z^{k-j}(1+z)^{i+j}.
\end{split}
\end{equation}
由于 $[z^k]z^{k-j}(1+z)^{i+j}=\binom{i+j}{j}$，得
\begin{equation}\label{eq:adiag}
a_k=\diag_k\Psi^k.
\end{equation}
定义
\begin{equation}\label{eq:c}
c_r=\diag_r\Phi^r\in\Z\quad(r\ge0),\qquad c_{-1}=0.
\end{equation}
由 $\Psi=\Phi+uvz$ 的二项式展开，
\begin{equation}\label{eq:bintransform}
a_k=\sum_{r=0}^k\binom kr c_r.
\end{equation}
再令
\begin{equation}\label{eq:b}
b_m=\frac{2(c_{m-1}+3c_{m-2})}{m}\quad(m\ge1);
\qquad b_1=2.
\end{equation}
利用 $k\binom kr=r\binom kr+(r+1)\binom{k}{r+1}$ 以及逐列求和，
\begin{equation}
\sum_{k=r}^{n-1}(3k+2)\binom kr
=(3r+2)\binom n{r+1}+3(r+1)\binom n{r+2}.
\end{equation}
代入 \eqref{eq:bintransform}，整理得
\begin{equation}
T_n=2\sum_{r=0}^{n-1}\binom{n-1}{r}
\frac{(3r+2)c_r+3rc_{r-1}}{r+1}.
\end{equation}
因为
\[
(3r+2)c_r+3rc_{r-1}
=3(r+1)(c_r+c_{r-1})-(c_r+3c_{r-1}),
\]
所以
\begin{equation}\label{eq:bridge}
\displaystyle
T_n=\sum_{r=0}^{n-1}\binom{n-1}{r}
\bigl(6(c_r+c_{r-1})-b_{r+1}\bigr).
\end{equation}
因此只需证明 $24b_m\in\Z$，即可得到 $24T_n\in\Z$。

\subsection{多项式恒等式与对角系数}

为消去式\eqref{eq:b}中的分母 $m$，将其分子表示为 Euler 导数的对角系数。
对角系数提取使 Euler 导数产生因子 $m$。令
\begin{equation}\label{eq:L}
\Lambda=(u+v+2uv)(u+v+3uv),
\end{equation}
并记 $D_u=u\partial_u$、$D_v=v\partial_v$、$D_z=z\partial_z$。

\begin{lemma}[多项式微分恒等式]\label{lem:certificate}
附录~\ref{sec:certificate}~给出的整数多项式 $P_u,P_v,P_z$ 满足
\begin{equation}\label{eq:certificate}
\Phi\Lambda\sum_iD_iP_i-\sum_iP_iD_i(\Phi\Lambda)
=48uvz(\Phi+3uvz)\Lambda^2,
\end{equation}
其中各求和遍历 $i\in\{u,v,z\}$。
\end{lemma}
\begin{proof}
将式\eqref{eq:polysP}代入，两端按 $z$ 展开后具有式\eqref{eq:certificate-table}所列的相同系数。
\end{proof}

多项式 $P_u,P_v,P_z$ 可用待定系数法求得。取分变量次数范围
\begin{equation}\label{eq:degree-box}
\deg_{u,v,z}P_u\le(4,3,2),\quad
\deg_{u,v,z}P_v\le(3,4,2),\quad
\deg_{u,v,z}P_z\le(3,3,2).
\end{equation}
三个次数上界依次对应变量 $u,v,z$。
比较式\eqref{eq:certificate}中各单项式的系数，得到关于168个待定系数的线性方程组。
附录~\ref{sec:certificate}~中的多项式给出其中一组整数解。

由于 $\Lambda(u,0)=u^2$ 是 $\Z[u,u^{-1}]$ 中的单位，$1/\Lambda$ 有如下关于 $v$ 的展开：

\begin{equation}\label{eq:Lexpansion}
\frac1\Lambda=\sum_{h,j\ge0}(-1)^{h+j}
v^{h+j}u^{-h-j-2}(1+2u)^h(1+3u)^j.
\end{equation}
在 $\Z[u,u^{-1}][[v]][z]$ 中，Euler 微分满足
\begin{equation}\label{eq:eulerdiag}
\diag_mD_iF=m\diag_mF\qquad(i=u,v,z).
\end{equation}
这是因为 $D_u,D_v,D_z$ 分别将单项式乘以相应变量的指数。

对 $m\ge2$，将 \eqref{eq:certificate} 乘以 $\Phi^{m-2}/(24\Lambda^2)$，由乘积求导法则得
\begin{equation}\label{eq:parts}
\begin{split}
2uvz(\Phi+3uvz)\Phi^{m-2}
=\frac1{24}\sum_i\left[
D_i\!\left(\frac{P_i\Phi^{m-1}}\Lambda\right)
-m\frac{P_i\Phi^{m-2}D_i\Phi}{\Lambda}\right].
\end{split}
\end{equation}
取 $\diag_m$，并应用 \eqref{eq:eulerdiag}，左侧为
$2(c_{m-1}+3c_{m-2})$，故
\begin{equation}\label{eq:bcertificate}
\displaystyle
b_m=\frac1{24}\diag_m\left[
\frac{\Phi^{m-2}}\Lambda\sum_iP_i(\Phi-D_i\Phi)\right].
\end{equation}
提取 $v^mz^m$ 的系数时，式\eqref{eq:Lexpansion}中只有 $h+j\le m$ 的项有贡献。
所得表达式是整数系数的 $u$-Laurent 多项式，故
\begin{equation}\label{eq:b24}
24b_m\in\Z\quad(m\ge2).
\end{equation}
$b_1=2$ 也满足此结论，由 \eqref{eq:bridge} 得
\begin{equation}\label{eq:T24}
24T_n\in\Z\qquad(n\ge1).
\end{equation}

\section{二进整性}
\begin{lemma}\label{lem:period2}
对所有整数 $k,s\ge0$，
\begin{equation}\label{eq:period2}
a_{k+2^s}\equiv a_k\pmod{2^s}.
\end{equation}
\end{lemma}
\begin{proof}
$s=0$ 时结论成立。设 $s\ge1$，在 $\Z_{(2)}$ 中写成
\begin{equation}\label{eq:a2}
a_k=9^k\sum_{i,j\ge0}\binom ki^2\binom kj^2
\binom{i+j}{i}\left(\frac23\right)^{i+j},
\end{equation}
其中 $i>k$ 时约定 $\binom ki=0$。模 $2^s$ 时只需保留 $i+j<s$ 的项。

对 $1\le i<s$，Vandermonde 恒等式给出
\[
\binom{k+2^s}{i}-\binom ki
=\sum_{h=1}^i\binom{2^s}{h}\binom{k}{i-h}.
\]
由 $h\binom{2^s}{h}=2^s\binom{2^s-1}{h-1}$，
\[
v_2\!\left(\binom{2^s}{h}\right)
\ge s-v_2(h)\ge s-\lfloor\log_2 i\rfloor.
\]
因此，平方之差满足
\begin{equation}\label{eq:2bindiff}
v_2\!\left(\binom{k+2^s}{i}^{\!2}-\binom ki^2\right)
\ge s-\lfloor\log_2 i\rfloor.
\end{equation}
乘上 \eqref{eq:a2} 中的 $2^{i+j}$ 后，该差被 $2^s$ 整除，因为
$i+j\ge\lfloor\log_2 i\rfloor$。$i=0$ 时因子不变；$j$ 因子同理。
再由反复平方可得 $9^{2^s}\equiv1\pmod{2^{s+3}}$，故 \eqref{eq:period2} 成立。
\end{proof}

令 $f(k)=(3k+2)a_k$，则
\begin{equation}\label{eq:fperiod}
f(k+2^s)\equiv f(k)\pmod{2^s}.
\end{equation}
对任意 $M\ge0,e\ge1$，用归纳法证明
\begin{equation}\label{eq:blocks}
2^{e-1}\ \bigg|\ \sum_{k=M}^{M+2^e-1}f(k).
\end{equation}
$e=1$ 时成立。对 $e\ge2$，将该区间等分，记两半之和为 $U,V$。
由 \eqref{eq:fperiod}，$V\equiv U\pmod{2^{e-1}}$。
归纳假设给出 $2^{e-2}\mid U$，于是 $U+V\equiv2U\equiv0\pmod{2^{e-1}}$。

若 $n=2^eq$，$q$ 为奇数且 $e\ge1$，将 $[0,n)$ 分成 $q$ 个长度为 $2^e$ 的区间，
由 \eqref{eq:blocks} 得 $2^{e-1}\mid\sum_{k<n}(3k+2)a_k$。故
\begin{equation}\label{eq:T2}
v_2(T_n)\ge0\qquad(n\ge1).
\end{equation}
$n$ 为奇数时此式由 \eqref{eq:T} 直接成立。

\section{三进整性}\label{sec:three-finite}
\subsection{奇次幂组合的局部整性引理}
置 $w=(1-16x)^{1/2}\in\Q[[x]]$，其中 $w(0)=1$。
对整数 $\alpha$ 和非负整数 $j$，有
\begin{equation}
\binom{\alpha/2}{j}\in\mathbb Z[1/2]
\qquad(\alpha\in\mathbb Z,\ j\ge0).
\label{eq:three-half-binomial}
\end{equation}
其分子和分母分别为
\(\alpha(\alpha-2)\cdots(\alpha-2j+2)\) 与 \(2^j j!\)。对每个奇素数幂 \(p^e\)，
因为 \(2\) 在模 \(p^e\) 下可逆，上述 \(j\) 个分子因子中至少有
\(\lfloor j/p^e\rfloor\) 个被 \(p^e\) 整除。这消去了 \(j!\) 的全部
奇素因子，证明了 \eqref{eq:three-half-binomial}；若分子为零则结论也成立。
特别地，每个整数次幂 \(w^\alpha\) 的所有形式幂级数系数都属于
\(\mathbb Z[1/2]\)。

对 \(H(x)=\sum_{k\ge0}h_kx^k\in\mathbb Q[[x]]\)，定义
\[
J_H(x)=\sum_{n\ge1}\frac{2}{n}
\left(\sum_{k=0}^{n-1}(3k+2)h_k\right)x^n.
\]
等价地，$J_H$ 由下列条件确定：
\begin{equation}
J_H(0)=0,
\qquad
J_H'(x)=\frac{2(3xH'(x)+2H(x))}{1-x}.
\label{eq:three-primitive-definition}
\end{equation}

\begin{lemma}[奇次幂组合的三进整性]\label{lem:odd-integrality}
设 $m\ge1$，且
\[
H(x)=\sum_{\substack{\alpha\in2\Z+1\\-4m-1\le\alpha\le m-1}}
\gamma_\alpha w^\alpha,\qquad \gamma_\alpha\in\Z_{(3)}.
\]
则 $3^mJ_H(x)\in\Z_{(3)}[[x]]$。
此外，对 $H(x)=w^{-1}$，有 $J_H(x)\in\Z_{(3)}[[x]]$。
\end{lemma}
\begin{proof}
由线性性，只需考虑单个奇次幂。
设 \(\alpha\) 是奇整数，并简记 \(J_\alpha=J_{w^\alpha}\)。由直接微分，
\begin{equation}
J_\alpha'(x)=
\frac{(3\alpha +4)w^\alpha-3\alpha w^{\alpha-2}}{1-x}.
\label{eq:three-primitive-derivative}
\end{equation}
对任意非负整数 \(N\)，令
\begin{equation}
\begin{aligned}
K_{\alpha,N}(x)=-2\sum_{\ell=0}^{N-1}(-15)^\ell\Biggl\{&
\frac{3\alpha +4}{\alpha-2\ell}\bigl(w^{\alpha-2\ell}-1\bigr)\\
&-\frac{3\alpha }{\alpha-2\ell-2}\bigl(w^{\alpha-2\ell-2}-1\bigr)
\Biggr\}.
\end{aligned}
\label{eq:three-finite-primitive}
\end{equation}
当 \(N=0\) 时，约定 \(K_{\alpha,0}=0\)。由于 \(\alpha\) 为奇数，式中分母均不为零。

有限几何恒等式给出
\begin{equation}
\frac{1}{w^2+15}
=\sum_{\ell=0}^{N-1}(-15)^\ell w^{-2\ell-2}
+\frac{(-15)^N w^{-2N}}{w^2+15}.
\label{eq:three-finite-geometric}
\end{equation}
利用 \(w^2+15=16(1-x)\) 并微分
\eqref{eq:three-finite-primitive}，得
\begin{equation}
J_\alpha'(x)-K_{\alpha,N}'(x)
=(-15)^N
\frac{(3\alpha +4)w^{\alpha-2N}-3\alpha w^{\alpha-2N-2}}{1-x}.
\label{eq:three-remainder-derivative}
\end{equation}
两边的原函数常数项均取零。因此，对每个正整数 \(n\)，
\begin{equation}
[x^n](J_\alpha-K_{\alpha,N})
=\frac{(-15)^N}{n}\,[x^{n-1}]
\frac{(3\alpha +4)w^{\alpha-2N}-3\alpha w^{\alpha-2N-2}}{1-x}.
\label{eq:three-remainder-coefficient}
\end{equation}
由式\eqref{eq:three-half-binomial}，式\eqref{eq:three-remainder-coefficient}中
\([x^{n-1}]\) 所提取的系数属于 \(\mathbb Z[1/2]\)，故
\begin{equation}
v_3\bigl([x^n](J_\alpha-K_{\alpha,N})\bigr)\ge N-v_3(n).
\label{eq:three-remainder-valuation}
\end{equation}

先设 \(m\ge1\)，并设奇整数 \(\alpha\) 满足
\(-4m-1\le \alpha\le m-1\)。对于
\eqref{eq:three-finite-primitive} 中的任一分母 \(d\)，有
\begin{equation}
0<|d|\le4m+3+2\ell<3^{m+\ell+1}.
\label{eq:three-denominator-bound}
\end{equation}
最后的不等式在 \((m,\ell)=(1,0)\) 时为 \(7<9\)。每次将 \(m\) 或
\(\ell\) 增加 \(1\)，左边分别增加 \(4\) 或 \(2\)，而右边乘以
\(3\)，故该不等式对所有 \(m\ge1,\ell\ge0\) 成立。
因此 \(v_3(d)\le m+\ell\)。式
\eqref{eq:three-finite-primitive} 中的因子 \(3^m15^\ell\) 消去每个
分母中的全部 \(3\) 因子。再由式\eqref{eq:three-half-binomial}，对每个整数 \(\beta\)，
\(w^\beta-1\) 的系数属于 \(\mathbb Z[1/2]\)，于是
\begin{equation}
v_3\bigl(3^m[x^n]K_{\alpha,N}\bigr)\ge0
\quad\text{对每个 }n\ge0\text{ 和 }N\ge0.
\label{eq:three-finite-main-integral}
\end{equation}

当 \(m=0,\alpha=-1\) 时，第一类分母的绝对值为
\(2\ell+1<3^{\ell+1}\)，其 \(3\) 阶至多为 \(\ell\)；第二类分母的
绝对值为 \(2\ell+3<3^{\ell+2}\)，其 \(3\) 阶至多为 \(\ell+1\)。
第二类分子 \(3\alpha=-3\) 含因子3，因此式\eqref{eq:three-finite-main-integral}仍成立。

固定正整数 \(n\)，取 \(N=v_3(n)\)。由式
\eqref{eq:three-remainder-valuation}，得
\[
v_3\bigl(3^m[x^n](J_\alpha-K_{\alpha,N})\bigr)\ge m\ge0.
\]
结合式\eqref{eq:three-finite-main-integral}，得到
\begin{equation}
v_3\bigl(3^m[x^n]J_\alpha\bigr)\ge0.
\label{eq:three-power-integral}
\end{equation}
此式对 $m\ge1$ 时指定范围内的奇整数 $\alpha$ 以及 $m=0,\alpha=-1$ 均成立。
再由 $J_H$ 的线性性和 $\gamma_\alpha\in\Z_{(3)}$，即得结论。
\end{proof}

\subsection{原序列的重新编号与奇次幂展开}

对非负整数 \(r,s\)，记
\[
m=r+s,\qquad R=\max(r,s),\qquad S=\min(r,s),
\]
并定义
\[
h_k^{(r,s)}=
\begin{cases}
\displaystyle
\binom{k}{r}^{\!2}\binom{k}{s}^{\!2}
\binom{2k-m}{k-r}4^k,&k\ge R,\\[2mm]
0,&k<R.
\end{cases}
\]
在 \(a_k\) 的定义中令 \(i=k-r\)、\(j=k-s\)，得
\begin{equation}
a_k=\sum_{r,s\ge0}3^{r+s}2^{-(r+s)}h_k^{(r,s)}.
\label{eq:three-reindex}
\end{equation}
对于每个 \(k\)，右边仅有 \(r,s\le k\) 的项可能非零。

定义
\[
F_{r,s}(x)=\sum_{k\ge R}\binom{2k-m}{k-r}4^kx^k,
\qquad
H_{r,s}(x)=\sum_{k\ge0}h_k^{(r,s)}x^k.
\]
有
\begin{equation}
F_{r,s}(x)=2^{-m}\frac{(1-w)^R(1+w)^S}{w}.
\label{eq:three-base-series}
\end{equation}

先证生成函数恒等式
\begin{equation}
\sum_{\ell\ge0}\binom{2\ell+d}{\ell}y^\ell
=\frac{C(y)^d}{\sqrt{1-4y}},
\qquad
C(y)=\frac{1-\sqrt{1-4y}}{2y},\quad d\ge0.
\label{eq:three-binomial-series}
\end{equation}
这里 \(C(0)=1\)，且 \(C-1=yC^2\)。为证明该恒等式，记其左边为
\(B_d(y)\)。二项式定理给出
\[
B_0(y)=\frac{1}{\sqrt{1-4y}},
\qquad
B_1(y)=\frac{B_0(y)-1}{2y}
=\frac{C(y)}{\sqrt{1-4y}}.
\]
其中第二式也直接来自
\(2\binom{2\ell+1}{\ell}=\binom{2\ell+2}{\ell+1}\)。
帕斯卡恒等式给出
\[
B_d=B_{d-1}+yB_{d+1}\qquad(d\ge1).
\]
对非负整数 $d$，令 $\widetilde B_d=C(y)^d/\sqrt{1-4y}$。
由 $C-1=yC^2$，序列 $\widetilde B_d$ 满足上述递推，且
$\widetilde B_0=B_0$、$\widetilde B_1=B_1$。
若 $d\ge1$，并已知 $B_{d-1}=\widetilde B_{d-1}$ 和 $B_d=\widetilde B_d$，
比较两序列的递推式，得
\[
y\bigl(B_{d+1}-\widetilde B_{d+1}\bigr)=0.
\]
由于 $\Q[[y]]$ 没有零因子，且 $y\ne0$，可约去 $y$，得到
$B_{d+1}=\widetilde B_{d+1}$。从 $d=1$ 起归纳，即得式\eqref{eq:three-binomial-series}。

由对称性可设 \(r\ge s\)。令 \(k=r+\ell\)、\(d=r-s\)，则
\[
F_{r,s}(x)=\frac{(4x)^r C(4x)^{r-s}}{w}.
\]
利用 \(4x=(1-w^2)/4\)、\(C(4x)=2/(1+w)\)，即得式\eqref{eq:three-base-series}。

\paragraph{二项式微分算子。}

令 \(\theta=x\,d/dx\)，并定义
\(\binom{\theta}{j}=\theta(\theta-1)\cdots(\theta-j+1)/j!\)，
其中 \(j=0\) 时取恒等算子。由 \(\binom{\theta}{j}x^k=\binom{k}{j}x^k\)，
有
\begin{equation}
H_{r,s}=\binom{\theta}{r}^{\!2}\binom{\theta}{s}^{\!2}F_{r,s}.
\label{eq:three-operators}
\end{equation}
此处的平方表示相应微分算子的复合。对任意整数 \(\alpha\) 和非负整数 \(j\)，
\begin{equation}
\binom{\theta}{j}w^\alpha
=\frac{x^j}{j!}\frac{d^j}{dx^j}w^\alpha
=\binom{\alpha/2}{j}(w^2-1)^j w^{\alpha-2j}.
\label{eq:three-operator-power}
\end{equation}

由 \eqref{eq:three-base-series}，\(F_{r,s}\) 是系数属于
\(\mathbb Z[1/2]\)、指数介于 \(-1\) 与 \(m-1\) 之间的 \(w\)-Laurent
多项式。公式 \eqref{eq:three-operator-power} 表明，每个算子
\(\binom{\theta}{j}\) 保持这一系数环、保持指数的奇偶性，并使指数
至多降低 \(2j\)。式 \eqref{eq:three-operators} 中四个算子的下标之和
为 \(2m\)，故 \(H_{r,s}\) 的指数介于 \(-4m-1\) 与 \(m-1\) 之间。

如果 \(m>0\)，\(F_{r,s}\) 的偶数次幂部分是 \(w^2\) 的多项式，因此
也是 \(x\) 的多项式，次数至多为
\(\lfloor(m-1)/2\rfloor<R\)。而式 \eqref{eq:three-operators} 的算子
消去每个次数小于 \(R\) 的单项式 \(x^k\)，因为此时
\(\binom{k}{r}\binom{k}{s}=0\)。故偶数次幂部分为零。
当 \(m=0\) 时，\(H_{0,0}=w^{-1}\)。因此
\begin{equation}
H_{r,s}(x)=
\sum_{\substack{\alpha\in2\Z+1\\-4m-1\le\alpha\le m-1}}
\gamma_\alpha w^\alpha,
\qquad \gamma_\alpha\in\mathbb Z[1/2].
\label{eq:three-odd-expansion}
\end{equation}
由定义，\(H_{r,s}\) 中次数小于 \(R\) 的系数均为零。

\subsection{加权和的三进整性}

对有限展开 \eqref{eq:three-odd-expansion} 应用引理~\ref{lem:odd-integrality}，利用 \(J_H\) 对 \(H\) 的线性，得到
\[
v_3\left(3^{r+s}2^{-(r+s)}[x^n]J_{H_{r,s}}\right)\ge0
\qquad(r,s\ge0,\ n\ge1).
\]
因为 \(H_{r,s}\) 从次数 \(R\) 起，按
\eqref{eq:three-primitive-definition}，\(J_{H_{r,s}}\) 从次数
\(R+1\) 起。因此，对于固定的 \(n\)，只有有限多个
\(0\le r,s<n\) 的项可能贡献 \(x^n\) 的系数。
由式\eqref{eq:three-reindex}，得
\[
T_n=
\sum_{0\le r,s<n}3^{r+s}2^{-(r+s)}[x^n]J_{H_{r,s}}.
\]
每个加数的 \(3\) 阶均非负，所以
\begin{equation}\label{eq:T3}
v_3(T_n)\ge0\qquad(n\ge1).
\end{equation}

\section{主定理的证明}
\begin{proof}[定理~\ref{thm:main}~的证明]
由式\eqref{eq:T24}，$T_n$ 的最简分母只能含素因子2和3。
式\eqref{eq:T2}和\eqref{eq:T3}分别给出 $T_n\in\Z_{(2)}$ 和 $T_n\in\Z_{(3)}$。
因此 $T_n$ 的最简分母为1，即 $T_n\in\Z$，这与式\eqref{eq:goal}等价。
当 $n=1$ 时，$a_0=1$、$T_1=4$，亦符合结论。
\end{proof}

\appendix
\section{多项式证书}\label{sec:certificate}
定义以下五个二元整数多项式：
\begin{equation}\label{eq:polyA}
\begin{split}
A(u,v)={}&(2u+1)(3u+1)\bigl[
6(2u+1)(3u+1)v^3\\
&\quad +(42u^2+23u+1)v^2+5u(4u+1)v+4u^2\bigr].
\end{split}
\end{equation}
\begin{equation}\label{eq:polyB}
\begin{split}
B(u,v)={}&18(u+1)(2u+1)(3u+1)(6u+1)v^3\\
&+(792u^4+1464u^3+856u^2+184u+11)v^2\\
&+u(318u^3+484u^2+209u+22)v\\
&+u^2(2u+1)(21u+11).
\end{split}
\end{equation}
\begin{equation}\label{eq:polyC}
\begin{split}
C(u,v)={}&(2u+1)\bigl[3(36u^2+29u+7)v^3\\
&\quad +(132u^2+77u+10)v^2+u(53u+17)v+7u^2\bigr].
\end{split}
\end{equation}
\begin{align}
D(u,v)&=uv\bigl[4(12u^2+9u+1)v+u(51u+28)\bigr],\label{eq:polyD}\\
E(u,v)&=3uv(u-v)(56uv+25u+25v+8).\label{eq:polyE}
\end{align}
令
\begin{equation}\label{eq:polysP}
\begin{aligned}
P_u&=-A(u,v)-B(u,v)z-C(u,v)z^2,\\
P_v&=-C(v,u)-B(v,u)z-A(v,u)z^2,\\
P_z&=-D(u,v)-E(u,v)z+D(v,u)z^2.
\end{aligned}
\end{equation}
为核验 \eqref{eq:certificate}，记其左端为
\[
\mathcal E=\Phi \Lambda\sum_iD_iP_i-\sum_iP_iD_i(\Phi \Lambda).
\]
该式的 $z$ 次数至多为 $4$。将上述多项式代入并按 $z$ 展开，得到
\begin{equation}\label{eq:certificate-table}
\begin{array}{c|l}
j&[z^j]\mathcal E/(48uv\Lambda^2)\\ \hline
0&0\\
1&(2u+1)(3u+1)(2v+1)\\
2&\begin{aligned}36u^2v^2+30u^2v+6u^2+30uv^2+31uv\\{}+7u+6v^2+7v+2\end{aligned}\\
3&(2u+1)(2v+1)(3v+1)\\
4&0
\end{array}
\end{equation}
另一方面，$\Phi$ 的常数项和二次项分别是表中的 $j=1,3$ 两项，而
\begin{equation}
\begin{split}
[z]\Phi+3uv
&=(2u+1)(2v+1)(2+3u+3v+9uv)+2uv\\
&=36u^2v^2+30u^2v+6u^2+30uv^2\\
&\quad+31uv+7u+6v^2+7v+2.
\end{split}
\end{equation}
因此 \eqref{eq:certificate-table} 正是
$48uvz(\Phi+3uvz)\Lambda^2$ 的全部 $z$ 系数，恒等式 \eqref{eq:certificate} 得证。

\section{对角系数的有限多项式表示}\label{app:finite-diagonal}
% A self-contained finite-polynomial appendix; Phi, \Lambda and P_i are the certificate polynomials.
记
\[
 K=\sum_{i\in\{u,v,z\}}P_i(\Phi-D_i\Phi),\qquad
 c_r=[u^rv^rz^r]\Phi^r,\qquad
 b_m=\frac2m(c_{m-1}+3c_{m-2}).
\]
式\eqref{eq:bcertificate}也可写成有限多项式的系数。仍约定 $c_{-1}=0$。
对 $q\ge0$ 定义
\[
 U_q(u)=\sum_{t=0}^q\binom{q+1}{t+1}
        (3^{t+1}-2^{t+1})u^t
       =\sum_{r=0}^q(1+2u)^r(1+3u)^{q-r},
\]
并对 $m\ge0$ 定义整数系数多项式
\[
 Q_m(u,v)=\sum_{q=0}^m(-1)^q u^{m-q}v^qU_q(u).
\]
则对于每个 $m\ge2$，有恒等式
\[
 \boxed{24b_m=[u^{2m+2}v^mz^m]Q_mK\Phi^{m-2}.}
\]
等价地，
\[
\boxed{
24b_m=\sum_{q=0}^m\sum_{t=0}^q(-1)^q\binom{q+1}{t+1}
 (3^{t+1}-2^{t+1})
 [u^{m+q+2-t}v^{m-q}z^m]K\Phi^{m-2}.
}
\]
右端各项均为整数。

\paragraph{证明。}
令 $A_0=1+2u$, $B_0=1+3u$。由 $U_q$ 的有限和表达式可得
$U_0=1$, $U_1=A_0+B_0$ 及
$U_q=(A_0+B_0)U_{q-1}-A_0B_0U_{q-2}$。
因为 $\Lambda=(u+A_0v)(u+B_0v)$，逐项相乘得
\[
 \Lambda Q_m=u^{m+2}+(-1)^mv^{m+1}
       \bigl(uU_{m+1}+A_0B_0vU_m\bigr).
\]
以下 $F\equiv G\pmod{v^{m+1}}$ 只表示两有限多项式的
$v^0,\ldots,v^m$ 系数相等。
固定 $m\ge2$，令 $a=m+2$；于是 $\Lambda Q_m\equiv u^a$，且
\[
 \Lambda D_iQ_m+Q_mD_i\Lambda\equiv a\,\delta_{i,u}\,u^a.
\]
其中 $i\in\{u,v,z\}$，$\delta_{i,u}$ 为 Kronecker 符号：
$i=u$ 时取1，$i=v$ 或 $i=z$ 时取0。
这是因为 $D_u(u^a)=au^a$，而 $D_v(u^a)=D_z(u^a)=0$。
Euler 微分 $D_u,D_v,D_z$ 都保持上述系数相等关系。
定义有限多项式
\[
\begin{split}
\mathcal E_m={}&\sum_iD_i\bigl(P_iQ_m\Phi^{m-1}\bigr)
  -aP_uQ_m\Phi^{m-1}\\
 &-mQ_m\Phi^{m-2}\sum_iP_iD_i\Phi.
\end{split}
\]
应用乘积法则、$\Lambda Q_m\equiv u^a$ 及其三个微分式，可得
\[
\begin{split}
\Lambda^2\mathcal E_m
&\equiv u^a\Phi^{m-2}
 \left(\Phi \Lambda\sum_iD_iP_i-\sum_iP_iD_i(\Phi \Lambda)\right)\\
&=48u^auvz(\Phi+3uvz)\Phi^{m-2}\Lambda^2.
\end{split}
\]
由于 $\Lambda(u,0)^2=u^4\ne0$，可以在上述系数相等关系中约去 $\Lambda^2$：
若一个非零多项式最低的 $v$ 次数不超过 $m$，乘以 $\Lambda^2$ 后最低的
$v$ 次数不变，因为该最低项的系数只是乘以非零多项式 $u^4$。
因此
\[
\mathcal E_m\equiv48u^auvz(\Phi+3uvz)\Phi^{m-2}.
\]
对两端取 $[u^{2m+2}v^mz^m]$。
Euler 微分作用于该系数时，$D_u,D_v,D_z$ 分别给出乘数
$2m+2,m,m$；又有 $2m+2-a=m$。故左端系数为
\[
m[u^{2m+2}v^mz^m]Q_m\Phi^{m-2}
 \sum_iP_i(\Phi-D_i\Phi).
\]
右端系数为
$48(c_{m-1}+3c_{m-2})=24m b_m$。
约去 $m$ 即得上述公式。
$m=1$ 时，$b_1=2$。

\section{计算核实}\label{app:reproducibility}
式\eqref{eq:certificate}及其系数表\eqref{eq:certificate-table}已用 Python 程序核实。
计算采用任意精度整数运算，逐项比较两端的全部系数，所得残差均为零。

\begin{thebibliography}{9}
\bibitem{Guo2016}
Guo, V.J.W., Proof of Sun's conjectures on integer-valued polynomials,
J. Math. Anal. Appl., 2016, 444(1): 182--191.
\href{https://doi.org/10.1016/j.jmaa.2016.06.028}{doi:10.1016/j.jmaa.2016.06.028}.
\bibitem{Sun2012}
Sun, Z.-W., On sums of Ap\'ery polynomials and related congruences,
J. Number Theory, 2012, 132(11): 2673--2699.
\href{https://arxiv.org/abs/1101.1946}{arXiv:1101.1946}.
\bibitem{Sun2026}
Sun, Z.-W., A family of polynomials and related congruences and series,
Front. Comb. Number Theory, 2026, 2: 1--22.
\href{https://doi.org/10.3934/fcnt.2026009}{doi:10.3934/fcnt.2026009}.
\bibitem{Zeilberger1991}
Zeilberger, D., The method of creative telescoping,
J. Symbolic Comput., 1991, 11(3): 195--204.
\href{https://doi.org/10.1016/S0747-7171(08)80044-2}{doi:10.1016/S0747-7171(08)80044-2}.
\end{thebibliography}
\end{document}
